% Vietnamese translation for OI Public Library
% Source-File: IOI中国国家候选队论文/2006/朱晨光/朱晨光.doc
\documentclass[11pt]{article}
\usepackage{oipl-vn}

\newcommand{\Oh}{\mathcal{O}}
\newcommand{\matzeroone}{ma trận \(0\)-\(1\)}

\title{Ứng dụng của các cấu trúc dữ liệu cơ bản trong thi tin học}
\author{Zhu Chenguang\\Trường THPT số 1 Vu Hồ, An Huy}
\date{Trại đông Tin học toàn quốc, 2006}

\begin{document}
\maketitle

\origtitle{Ứng dụng của các cấu trúc dữ liệu cơ bản trong thi tin học}
\sourcepath{IOI China candidate team papers / 2006 / Zhu Chenguang.doc}

\renewcommand{\abstractname}{Tóm tắt}
\begin{abstract}
Bài viết giới thiệu một số cấu trúc dữ liệu cơ bản, như danh sách tuyến tính,
ngăn xếp và hàng đợi, cùng cách chúng được dùng trong các bài toán thi tin học.
Thông qua ba ví dụ, tác giả nhấn mạnh rằng những cấu trúc tưởng như đơn giản
có thể thay thế các cấu trúc phức tạp hơn, giảm độ khó cài đặt mà vẫn giữ được
độ phức tạp tốt. Cuối bài viết là phần nhìn lại quan hệ giữa cấu trúc dữ liệu
cơ bản và cấu trúc dữ liệu nâng cao.
\end{abstract}

\noindent\textbf{Từ khóa:} cấu trúc dữ liệu cơ bản; danh sách tuyến tính;
hàng đợi; danh sách liên kết đôi; ngăn xếp; độ phức tạp cài đặt; độ phức tạp
thời gian; quan hệ biện chứng.

\section{Dẫn nhập}

Trong các kỳ thi tin học ngày nay, bài khó xuất hiện ngày càng nhiều. Đi cùng
với chúng thường là độ phức tạp cài đặt rất cao. Khoa học máy tính phát triển
đã sinh ra nhiều cấu trúc dữ liệu hiệu quả và thực dụng, nhưng độ phức tạp
cài đặt của chúng khiến thí sinh phải hết sức thận trọng; chỉ một sơ suất nhỏ
cũng có thể làm hỏng toàn bộ chương trình.

Tuy nhiên, không phải bài nào cũng cần cấu trúc dữ liệu phức tạp. Nhiều lúc,
những cấu trúc cơ bản mà ta dễ bỏ qua lại có đất dụng võ. Vận dụng linh hoạt
các cấu trúc ấy có thể tiết kiệm thời gian quý trong phòng thi và tăng xác
suất viết đúng chương trình.

\section{Một số cấu trúc dữ liệu cơ bản}

Phần này chỉ nhắc lại ngắn gọn các khái niệm vốn có trong mọi giáo trình cấu
trúc dữ liệu.

\subsection{Danh sách tuyến tính}

Danh sách tuyến tính là một trong các cấu trúc thường gặp và đơn giản nhất.
Nói ngắn gọn, một danh sách tuyến tính là một dãy hữu hạn gồm \(n\) phần tử dữ
liệu. Các phép toán cơ bản trên danh sách tuyến tính gồm:
\begin{enumerate}
  \item \texttt{INITIATE(L)}: khởi tạo;
  \item \texttt{LENGTH(L)}: lấy độ dài;
  \item \texttt{GET(L,i)}: lấy phần tử thứ \(i\);
  \item \texttt{PRIOR(L,element)}: lấy phần tử đứng trước;
  \item \texttt{NEXT(L,element)}: lấy phần tử đứng sau;
  \item \texttt{LOCATE(L,x)}: định vị phần tử;
  \item \texttt{INSERT(L,i,b)}: chèn trước vị trí \(i\);
  \item \texttt{DELETE(L,i)}: xóa vị trí \(i\);
  \item \texttt{EMPTY(L)}: kiểm tra rỗng;
  \item \texttt{CLEAR(L)}: xóa toàn bộ danh sách.
\end{enumerate}

\paragraph{Lưu trữ tuần tự.}
Trong máy tính, cách biểu diễn đơn giản và phổ biến nhất là dùng một dãy ô nhớ
có địa chỉ liên tiếp để lưu lần lượt các phần tử; nói cách khác, đó là mảng
một chiều. Nếu phần tử đầu ở địa chỉ \(b\), mỗi phần tử chiếm \(L\) đơn vị bộ
nhớ, thì phần tử thứ \(i\) ở địa chỉ \(b+(i-1)L\). Cách lưu trữ này cho phép
truy cập ngẫu nhiên: lấy phần tử bất kỳ trong \(O(1)\). Trong Pascal hoặc C,
ta mô tả nó bằng mảng một chiều.

Trong cấu trúc lưu trữ tuần tự, chèn hoặc xóa một phần tử bất kỳ tốn
\(\Oh(n)\), vì có thể phải dời nhiều phần tử; định vị theo chỉ số tốn
\(\Oh(1)\).

\paragraph{Lưu trữ liên kết.}
Với danh sách liên kết đơn, mỗi nút lưu dữ liệu và một con trỏ tới nút kế
tiếp. Nút đầu giả \texttt{head} giữ con trỏ tới phần tử đầu tiên. Các nút
không cần nằm liên tiếp trong bộ nhớ.

Danh sách liên kết vòng là biến thể trong đó con trỏ của nút cuối trỏ lại nút
đầu giả, tạo thành một vòng. Danh sách liên kết đôi thì mỗi nút có hai con
trỏ: một tới nút kế tiếp và một tới nút đứng trước. Nó cũng có thể được tổ
chức thành dạng vòng.

Trong cấu trúc lưu trữ liên kết, nếu vị trí đã được định vị thì chèn hoặc xóa
tốn \(\Oh(1)\). Ngược lại, tìm vị trí của một phần tử thường tốn \(\Oh(n)\).

\subsection{Ngăn xếp}

Ngăn xếp là danh sách tuyến tính chỉ cho phép chèn và xóa ở cuối danh sách.
Đầu cuối ấy gọi là đỉnh ngăn xếp; đầu còn lại gọi là đáy. Ngăn xếp rỗng là
ngăn xếp không có phần tử.

Các thao tác sửa đổi của ngăn xếp tuân theo nguyên tắc vào sau ra trước
(\term{last in, first out}, viết tắt là LIFO). Các phép toán cơ bản gồm:
\begin{itemize}
  \item \texttt{INISTACK(S)}: khởi tạo;
  \item \texttt{EMPTY(S)}: kiểm tra rỗng;
  \item \texttt{PUSH(S,x)}: đưa \(x\) vào ngăn xếp;
  \item \texttt{POP(S,x)}: lấy phần tử ở đỉnh ra;
  \item \texttt{GETTOP(S)}: xem phần tử ở đỉnh;
  \item \texttt{CLEAR(S)}: xóa toàn bộ;
  \item \texttt{CURRENT\_SIZE(S)}: lấy số phần tử hiện có.
\end{itemize}
Mỗi thao tác cơ bản trên ngăn xếp đều tốn \(\Oh(1)\).

\subsection{Hàng đợi}

Ngược với ngăn xếp, hàng đợi là danh sách tuyến tính vào trước ra trước
(\term{first in, first out}, viết tắt là FIFO). Nó chỉ cho phép chèn ở một đầu
và xóa ở đầu còn lại. Đầu được chèn gọi là cuối hàng; đầu được xóa gọi là đầu
hàng.

Các phép toán cơ bản gồm:
\begin{itemize}
  \item \texttt{INIQUEUE(Q)}: khởi tạo;
  \item \texttt{EMPTY(Q)}: kiểm tra rỗng;
  \item \texttt{ENQUEUE(Q,x)}: đưa \(x\) vào hàng;
  \item \texttt{DLQUEUE(Q)}: lấy phần tử ở đầu hàng ra;
  \item \texttt{GETHEAD(Q)}: xem phần tử ở đầu hàng;
  \item \texttt{CLEAR(Q)}: xóa toàn bộ;
  \item \texttt{CURRENT\_SIZE(Q)}: lấy số phần tử hiện có.
\end{itemize}
Mỗi thao tác cơ bản trên hàng đợi đều tốn \(\Oh(1)\).

Ngoài ngăn xếp và hàng đợi còn có hàng đợi hai đầu. Ở đó, mỗi đầu có thể chỉ
nhập, chỉ xuất, hoặc vừa nhập vừa xuất. Hàng đợi thường có thể xem là một
trường hợp đặc biệt của hàng đợi hai đầu. Các cấu trúc cơ bản khác như chuỗi
và cách lưu cây có gốc cũng rất quan trọng, nhưng bài viết này tập trung vào
ngăn xếp, hàng đợi và danh sách tuyến tính.

\section{Ứng dụng của ngăn xếp}

\subsection{Hình chữ nhật toàn số không lớn nhất}

\paragraph{Bài toán.}
Cho một \matzeroone{} kích thước \(m\times n\), hãy tìm diện tích hình chữ
nhật lớn nhất chỉ gồm các ô bằng \(0\).

\paragraph{Phân tích.}
Cách trực tiếp là liệt kê góc trái trên và góc phải dưới, rồi kiểm tra từng ô
bên trong. Độ phức tạp là \(\Oh(m^3n^3)\). Dùng tổng tiền tố hai chiều có thể
kiểm tra một hình chữ nhật trong \(\Oh(1)\), giảm tổng thời gian xuống
\(\Oh(m^2n^2)\), nhưng vẫn chưa đủ tốt.

Trong luận văn đội tuyển quốc gia năm 2003, Wang Zhikun đã nghiên cứu bài này
và đưa ra thuật toán \(\Oh(mn)\) dựa trên quy hoạch động. Bài viết này dùng
một tính chất then chốt trong cách nhìn đó: trong mọi hình chữ nhật toàn số
không cực đại đều tồn tại một ``đường treo'' thẳng đứng. Đường treo bắt đầu
từ một ô và kéo lên trên cho tới khi gặp số \(1\) hoặc biên ma trận; hình chữ
nhật cực đại có thể thu được bằng cách mở rộng đường treo ấy sang trái và
sang phải.

Gọi \(h(x,y)\) là độ dài đường treo kết thúc tại ô \((x,y)\). Ta có:
\[
h(x,y)=
\begin{cases}
0, & \text{nếu ô }(x,y)\text{ là }1,\\
1, & \text{nếu ô }(x,y)\text{ là }0\text{ và }x=1,\\
h(x-1,y)+1, & \text{nếu ô }(x,y)\text{ là }0\text{ và }x>1.
\end{cases}
\]

Sau khi tính \(h(x,y)\) cho một hàng \(x\), ta cần biết mỗi đường treo có thể
mở rộng sang trái và sang phải bao xa mà không gặp chướng ngại. Cách làm trong
bài này là duyệt hàng ấy bằng một ngăn xếp đơn điệu.

Để tránh xử lý trường hợp cuối hàng, thêm một lính canh \(h(x,n+1)=-1\). Trước
khi xử lý mỗi hàng, làm rỗng ngăn xếp. Mỗi phần tử trong ngăn xếp lưu hai giá
trị: chiều cao \(v\) và cột trái nhất \(j\) mà một đường treo có chiều cao
\(v\) có thể vươn tới.

Khi xét cột \(i\), gọi phần tử đỉnh là \((A,j)\).
\begin{itemize}
  \item Nếu ngăn xếp rỗng hoặc \(h(x,i)>A\), đưa \((h(x,i),i)\) vào ngăn xếp.
  \item Nếu \(h(x,i)=A\), không cần chèn thêm, vì đường treo mới tạo ra cùng
  một tập hình chữ nhật như phần tử đã có, còn cột \(j\) đã ghi vị trí trái
  nhất.
  \item Nếu \(h(x,i)<A\), liên tục lấy ra các phần tử cao hơn. Với mỗi phần
  tử \((B,j)\) bị lấy ra, nó đã gặp biên phải ở cột \(i\), nên diện tích ứng
  viên là \(B(i-j)\). Sau khi lấy ra xong, nếu đỉnh mới có chiều cao bằng
  \(h(x,i)\) thì không cần chèn; ngược lại chèn \(h(x,i)\) với cột trái nhất
  bằng cột \(j\) của phần tử cuối cùng vừa bị lấy ra.
\end{itemize}

Trong mô tả hình của bản gốc, các phần tử trong ngăn xếp được vẽ như các điểm
trên hệ trục: trục hoành là cột trái nhất, trục tung là chiều cao. Khi một
chiều cao nhỏ hơn xuất hiện, các điểm cao hơn ở bên phải bị xóa vì chúng không
thể tiếp tục mở rộng qua cột hiện tại.

\paragraph{Giả mã.}
\begin{verbatim}
maxarea = 0
for x = 1..m:
    compute h(x, i) for all columns i
    h(x, n + 1) = -1
    clear stack
    for i = 1..n + 1:
        if stack is empty or h(x, i) > top.height:
            push(height = h(x, i), left = i)
        else if h(x, i) == top.height:
            continue
        else:
            lastLeft = i
            while stack is not empty and h(x, i) < top.height:
                lastLeft = top.left
                maxarea = max(maxarea, top.height * (i - top.left))
                pop
            if stack is empty or h(x, i) > top.height:
                push(height = h(x, i), left = lastLeft)
\end{verbatim}

Mỗi phần tử được đưa vào và lấy ra khỏi ngăn xếp nhiều nhất một lần trong mỗi
hàng, nên tổng độ phức tạp là \(\Oh(mn)\). Dùng mảng cuộn cho \(h\) thì bộ nhớ
còn \(\Oh(n)\). Điểm khác với ngăn xếp thường là các chiều cao trong ngăn xếp
luôn tăng nghiêm ngặt từ đáy tới đỉnh, và mỗi phần tử còn mang theo vị trí
trái nhất tương ứng.

\section{Ứng dụng của danh sách tuyến tính}

\subsection{Thống kê doanh thu}

\paragraph{Bài toán.}
Cho doanh thu \(a_1,a_2,\ldots,a_n\) của \(n\) ngày, với
\(1\le n\le 32767\). Giá trị dao động nhỏ nhất của ngày \(i\) được định nghĩa
là
\[
  \min_{1\le j<i}|a_j-a_i|.
\]
Riêng ngày đầu tiên có giá trị dao động nhỏ nhất bằng \(a_1\). Hãy tính tổng
các giá trị dao động nhỏ nhất của mọi ngày.

\paragraph{Phân tích.}
Vòng lặp hai tầng cho độ phức tạp \(\Oh(n^2)\), không đủ nhanh. Nếu dùng cây
cân bằng, khi chèn doanh thu của ngày hiện tại ta có thể lấy giá trị gần nhất
trên đường tìm kiếm và đạt \(\Oh(n\log n)\). Nhưng cây cân bằng có độ phức tạp
cài đặt cao: quay trái, quay phải và các biến thể đều dễ sai trong điều kiện
thi.

Ta có thể đạt cùng bậc thời gian bằng danh sách liên kết đôi. Trước hết sắp
xếp \(n\) phần tử theo giá trị, đồng thời ghi vị trí \(b_i\) của phần tử
\(a_i\) trong dãy sau sắp xếp. Trên dãy đã sắp xếp, dựng một danh sách liên
kết đôi. Sau đó xử lý các phần tử theo thứ tự ngược từ \(a_n\) về \(a_1\).

Khi xử lý \(a_i\), các phần tử \(a_{i+1},a_{i+2},\ldots,a_n\) đã bị xóa khỏi
danh sách. Vì vậy, tiền nhiệm và hậu nhiệm của vị trí \(b_i\) trong danh sách
hiện tại chính là hai giá trị lân cận của \(a_i\) trong tập
\(\{a_1,\ldots,a_i\}\) sau khi sắp xếp. Dao động nhỏ nhất của ngày \(i\) chỉ
có thể là khoảng cách tới một trong hai giá trị đó. Sau khi cộng đáp án, xóa
\(a_i\) khỏi danh sách trong \(\Oh(1)\).

Ví dụ \(a_1=9,a_2=3,a_3=8\). Sau khi sắp xếp ta có \(3,8,9\). Với \(a_3=8\),
hai láng giềng là \(3\) và \(9\), nên dao động là \(1\). Xóa \(8\). Với
\(a_2=3\), chỉ còn hậu nhiệm \(9\), nên dao động là \(6\). Ngày đầu đóng góp
\(9\). Tổng là \(16\).

\paragraph{Giả mã.}
\begin{verbatim}
input n and a[1..n]
sort the indices by increasing a[index]
record pos[i] = position of a[i] in the sorted order
build a doubly linked list over the sorted indices

answer = a[1]
for i = n downto 2:
    best = infinity
    if previous[i] exists:
        best = min(best, abs(a[i] - a[previous[i]]))
    if next[i] exists:
        best = min(best, abs(a[i] - a[next[i]]))
    answer += best
    remove i from the doubly linked list
print answer
\end{verbatim}

Sắp xếp tốn \(\Oh(n\log n)\), còn dựng danh sách và xử lý đều tốn
\(\Oh(n)\). Vì vậy tổng độ phức tạp là \(\Oh(n\log n)\), ngang với cách dùng
cây cân bằng nhưng mã ngắn và ít rủi ro hơn.

\section{Ứng dụng của hàng đợi}

\subsection{Điệu valse lộng lẫy}

\paragraph{Bài toán.}
Cho một ma trận \(n\times m\), một số ô là chướng ngại. Một người bắt đầu ở
một ô trống và thực hiện \(k\) lượt trượt. Ở lượt thứ \(t\), hướng trượt đã
biết trước là đông, nam, tây hoặc bắc, và người đó có thể đi liên tiếp tối đa
\(c_t\) ô theo hướng ấy, cũng có thể đứng yên. Trong quá trình trượt không
được chạm chướng ngại. Hãy tìm quãng đường trượt dài nhất. Dữ liệu trong bản
gốc có \(1\le n,m\le 200\), \(k\le 200\), tổng thời gian không quá \(40000\).

\paragraph{Quy hoạch động trực tiếp.}
Đặt \(f(t,x,y)\) là quãng đường dài nhất sau \(t\) lượt nếu kết thúc ở ô
\((x,y)\). Với lượt trượt sang đông, ta có dạng chuyển:
\[
\begin{aligned}
f(t,x,y)=\max\{&
f(t-1,x,y),\\
&f(t-1,x,y-1)+1,\\
&f(t-1,x,y-2)+2,\ldots,\\
&f(t-1,x,y')+y-y'\},
\end{aligned}
\]
trong đó \(y'\) là cột trái nhất còn có thể đi tới mà không vượt quá \(c_t\)
ô và không đi xuyên qua chướng ngại. Ba hướng còn lại tương tự.

Số trạng thái là \(\Oh(knm)\), nhưng mỗi chuyển có thể phải xét
\(\Oh(\max\{n,m\})\) ô, nên độ phức tạp quá lớn.

\paragraph{Biến đổi công thức.}
Tạm bỏ qua chướng ngại. Với một hàng cố định và hướng sang đông, nếu
\(c_t=2\), ta có:
\[
\begin{array}{rcl}
f(t,x,1)&=&\max\{f(t-1,x,1)\},\\
f(t,x,2)&=&\max\{f(t-1,x,1)+1,\ f(t-1,x,2)\},\\
f(t,x,3)&=&\max\{f(t-1,x,1)+2,\ f(t-1,x,2)+1,\ f(t-1,x,3)\},\\
f(t,x,4)&=&\max\{f(t-1,x,2)+2,\ f(t-1,x,3)+1,\ f(t-1,x,4)\}.
\end{array}
\]
Đặt
\[
  a_i=f(t-1,x,i)-i+1.
\]
Khi đó các biểu thức trên trở thành việc lấy giá trị lớn nhất của một đoạn
liên tiếp trong dãy \(a\), rồi cộng thêm một hằng số theo cột hiện tại. Nếu có
chướng ngại, các đoạn bị cắt rời; khi gặp chướng ngại ta phải xóa toàn bộ cửa
sổ hiện tại.

Vì mỗi bước chỉ thêm một giá trị ở cuối cửa sổ và xóa một số giá trị ở đầu
cửa sổ, cấu trúc phù hợp tự nhiên là hàng đợi. Nhưng ta cần lấy giá trị lớn
nhất trong hàng, nên dùng \term{hàng đợi đơn điệu}.

\paragraph{Hàng đợi đơn điệu.}
Trong cửa sổ hiện tại, giả sử hai giá trị \(a_2\) và \(a_3\) cùng tồn tại và
\(a_2\le a_3\). Khi \(a_2\) chưa bị xóa, \(a_3\) chắc chắn cũng chưa bị xóa,
vì \(a_3\) được thêm sau. Do đó \(a_2\) không bao giờ có thể là giá trị lớn
nhất và không cần lưu. Từ nhận xét này suy ra các giá trị thật sự cần lưu
trong hàng đợi phải giảm nghiêm ngặt từ đầu tới cuối; phần tử đầu hàng chính
là giá trị lớn nhất của cửa sổ.

Khi xử lý một ô mới:
\begin{itemize}
  \item nếu ô là chướng ngại, xóa rỗng hàng đợi;
  \item xóa khỏi đầu hàng mọi phần tử đã cách ô hiện tại quá \(c_t\);
  \item trước khi chèn \(a_i\), xóa khỏi cuối hàng mọi phần tử không lớn hơn
  \(a_i\);
  \item chèn \(a_i\), rồi dùng phần tử đầu hàng để cập nhật \(f(t,x,y)\).
\end{itemize}

Hình minh họa của bản gốc vẽ các giá trị \(a_i\) trên hệ trục chỉ số--giá trị:
các điểm thấp hơn nằm trước một điểm cao hơn ở bên phải bị loại bỏ, vì chúng
không thể thắng trong bất kỳ cửa sổ tương lai nào.

\paragraph{Giả mã cho hướng đông.}
\begin{verbatim}
for x = 1..n:
    clear deque
    for y = 1..m:
        if cell (x, y) is blocked:
            f[now][x][y] = -infinity
            clear deque
            continue
        while deque is not empty and deque.front.index < y - c[t]:
            pop_front
        value = f[old][x][y] - y + 1
        while deque is not empty and deque.back.value <= value:
            pop_back
        push_back(index = y, value = value)
        f[now][x][y] = deque.front.value + y - 1
\end{verbatim}
Ba hướng còn lại chỉ khác thứ tự duyệt và cách tính khoảng cách. Mỗi ô vào và
ra hàng đợi nhiều nhất một lần trong mỗi lượt, vì vậy tổng độ phức tạp là
\(\Oh(knm)\), tốt hơn cách dùng cây đoạn \(\Oh(knm\log\max\{n,m\})\) và cũng
dễ cài đặt hơn.

\section{Tổng kết}

Ba ví dụ trên cho thấy các cấu trúc dữ liệu cơ bản không chỉ là kiến thức nhập
môn. Ngăn xếp, danh sách liên kết đôi và hàng đợi có thể, khi được tổ chức
đúng cách, giải những bài tưởng như cần cấu trúc dữ liệu cao cấp hơn.

Điểm mấu chốt không phải là phủ nhận cấu trúc dữ liệu nâng cao. Cây cân bằng,
cây đoạn và nhiều cấu trúc chuyên biệt vẫn cần thiết trong các tình huống phù
hợp. Nhưng trong thi đấu, ta luôn phải cân bằng giữa độ phức tạp thời gian,
độ phức tạp cài đặt và khả năng viết đúng. Cấu trúc cơ bản thường đem lại một
phương án có độ phức tạp đủ tốt với mã ngắn và ít lỗi.

Tác giả cũng nhấn mạnh một khía cạnh tư tưởng: việc quay lại dùng cấu trúc cơ
bản không phải là sự trở về đơn giản, càng không phải là loại bỏ tri thức nâng
cao. Ta đi từ cấu trúc cơ bản tới cấu trúc nâng cao, rồi nhìn lại cấu trúc cơ
bản ở một tầng hiểu biết cao hơn. Quá trình ấy giống một sự phát triển theo
đường xoắn ốc: bề ngoài là trở về điểm cũ, nhưng thực chất nhận thức đã sâu
hơn.

\section*{Tài liệu tham khảo}

\begin{enumerate}
  \item Yan Weimin, Wu Weimin. \textit{Data Structures}, ấn bản thứ hai.
  Thanh Hoa University Press, 1992.
  \item Thomas H. Cormen, Charles E. Leiserson, Ronald L. Rivest, Clifford
  Stein. \textit{Introduction to Algorithms}, Second Edition. The MIT Press,
  2001.
  \item Wang Zhikun. \textit{Bàn về cách dùng tư tưởng cực đại hóa để giải
  bài toán hình chữ nhật con lớn nhất}. Luận văn đội tuyển quốc gia IOI 2003.
\end{enumerate}

\section*{Lời cảm ơn}

Tác giả chân thành cảm ơn thầy Jiang Tao vì đã hướng dẫn và giúp đỡ trong quá
trình viết bài; đồng thời cảm ơn huấn luyện viên Liu Rujia và bạn Xu Zhilei
vì những góp ý quan trọng.

\section*{Phụ lục được lược dịch}

Bản gốc kèm các phát biểu đầy đủ của ba bài ví dụ và ba chương trình tham
khảo bằng C/C++. Theo yêu cầu của bản dịch này, phần chương trình dài không
được chép lại. Nội dung phụ lục có thể khái quát như sau:
\begin{itemize}
  \item Ví dụ 1 tương ứng bài \textit{Big Barn} của USACO Training: tìm hình
  vuông lớn nhất không chứa cây trên lưới \(N\times N\). Chương trình tham
  khảo dùng đúng ngăn xếp đơn điệu đã mô tả, chỉ thay diện tích hình chữ nhật
  bằng cạnh hình vuông lớn nhất.
  \item Ví dụ 2 là bài \textit{Thống kê doanh thu}: đọc \(n\) doanh thu, tính
  tổng dao động nhỏ nhất. Chương trình tham khảo sắp xếp chỉ số rồi xóa ngược
  trên danh sách liên kết đôi.
  \item Ví dụ 3 là bài \textit{Điệu valse lộng lẫy}, mã nguồn chính
  \texttt{adv1900}: dùng quy hoạch động hai lớp và hàng đợi đơn điệu theo từng
  hàng hoặc cột cho mỗi khoảng thời gian nghiêng của con tàu.
\end{itemize}

\end{document}
